\documentclass[aspectratio=169]{beamer}

\usepackage{slides/uqcs}

\title{Building with Bazel}
\author{Brae Webb}
\institute{University of Queensland}
\date{August 27, 2020}

\begin{document}

% Slide types
% title-slide
% split-slide
% section-slide
% body-slide
% image-slide
% 4-pane-slide

\begin{title-slide}
\titlepage
\end{title-slide}

\begin{split-slide}{Outline}
    \vspace{4em}
    \tableofcontents
\end{split-slide}

\section{What is Bazel?}

\begin{body-slide}{Bazel 101}
\begin{itemize}[<+-|alert@+>]
    \color{UQCSBlue}
    \item Pronounced same as `basil' (the herb) in US English `BAY-zel'
    \item Build tool for many many languages
    \item Open-source of Google's internal tool Blaze
\end{itemize}

\end{body-slide}

\section{An Ad for Bazel}

\subsection{Bazel Features}
\begin{split-slide}{Overview}
\begin{enumerate}
    \item Isolated
    \item Cached
    \item Dependency graph
    \item Extendable
\end{enumerate}
\end{split-slide}

\begin{body-slide}{Isolated}
Builds are isolated to enforce accurate dependency specification
\end{body-slide}

\begin{body-slide}{Cached}
Builds specify what they output and that is cached

The cache can be local or a remote cache
\end{body-slide}

\begin{body-slide}{Dependency Graph}
All dependencies are as accurate as possible

Dependency graph only updates what is required

Allows isolated builds to build in parallel
\end{body-slide}

\begin{body-slide}{Extendable}
Add new rules for a language or framework in a subset of python
\end{body-slide}

\subsection{Bazel Languages}
{
\usebackgroundtemplate{%
\includegraphics[width=\paperwidth,height=\paperheight]{images/languages}
}
\begin{frame}
    
\end{frame}
}


\section{Actually Using Bazel}

\subsection{BUILD}

\begin{body-slide}{BUILD}
\begin{itemize}[<+-|alert@+>] 
    \color{UQCSBlue}
    \item All the magic starts in a BUILD file
    \item Each 'package' in your repository will tend to have a BUILD
    \item BUILD files specify a list of build targets using build rules
\end{itemize}
\end{body-slide}

\begin{body-slide}{Built-in Rules}
Bazel comes with built-in rules that you can use straight away

\only<2->{
\begin{enumerate}
    \item C/C++ rules
    \item Java rules
    \item Python rules
    \item Protobuf rules
\end{enumerate}
}
\end{body-slide}

\begin{body-slide}{Example}
We have a Java program that serialises a tree
\end{body-slide}

\begin{body-slide}{Running Bazel}
\begin{itemize}[<+-|alert@+>] 
    \color{UQCSBlue}
    \item Bazel has 3 main modes
    \item \texttt{bazel build}
    \item \texttt{bazel run}
    \item \texttt{bazel test}
\end{itemize}
\end{body-slide}

\begin{body-slide}{Target Names}
\begin{itemize}[<+-|alert@+>] 
    \color{UQCSBlue}
    \item Relative path if in same directory
    \item Absolute path from anywhere in the repository
\end{itemize}
\end{body-slide}

\begin{body-slide}{Python Rules}
\begin{itemize}
    \item py\_library
    \item py\_binary
    \item py\_test
\end{itemize}
\end{body-slide}

\begin{body-slide}{Example}
We have a python file that computes anagrams
\end{body-slide}


\begin{body-slide}{Queries}
Bazel has a built-in query language to analyse your builds
\end{body-slide}

\subsection{WORKSPACE}

\begin{body-slide}{WORKSPACE}
\begin{itemize}[<+-|alert@+>]
    \color{UQCSBlue}
    \item Lives in the root of your repository
    \item Specifies external project dependencies
    \item Uses repository rules rather than rules
\end{itemize}
\end{body-slide}


\section{Extending Bazel}

\begin{body-slide}{Starlark}
Starlark is a language that is a subset of python

New rules are defined using the Starlark language
\end{body-slide}

\begin{body-slide}{Example}
Standard ML is a language for which we want to add support
\end{body-slide}

\begin{body-slide}{Further}
Sources for talk: \url{https://github.com/BraeWebb/bazel-talk}

Bazel documentation: \url{https://docs.bazel.build}

List of all things Bazel: \url{https://awesomebazel.com}
\end{body-slide}

\end{document}